%====================================================
% 第1章：緒論
%====================================================


\section{はじめに}
\subsection{研究背景}
近年，ロシアによるウクライナ侵略や中東地域における軍事的な緊張が高まり，化石燃料の調達に関する不確実性が
上昇するなど，日本が抱えるエネルギー需給構造上の課題が浮き彫りとなっている．
このような状況下において，省エネルギーに加え，再生可能エネルギー，原子力などの脱炭素電源を最大限活用する
ことが重要となっている．
特に，太陽光発電は比較的リードタイムが短く，分散型エネルギーリソースとしての活用が期待されているため，
更なる導入拡大が見込まれている\cite{ref01}．
再生可能エネルギーは変動が大きく,その変動に合わせて最適な電力変換が必要である．
再生可能エネルギーの普及拡大に伴い，系統を安定化させることが求められている\cite{ref02}．
\subsection{研究目的}
太陽光発電の電力源は直流出力のため，系統電圧に接続するためにはインバータが必要となる．
インバータ側の交流電圧と系統電圧の位相を同期するためには系統電圧の位相を検出する必要がある．
位相検出方法としてPLL(Phase Locked Loop)制御が用いられる．
系統電圧が三相である場合，三相二相変換，$dq$変換を用いて瞬時的な位相，振幅を検出することができる．
しかし，対象となる系統が単相である場合，系統電圧を単純に直交座標系に変換することができない為，
瞬時に位相角を検出することが困難である．
当研究室では，三角関数の積差の公式を用いて直交座標変換を行い，位相，振幅を
算出する疑似dq変換を用いたPLL制御を提案している．
疑似dq変換はサンプリング周波数が高くなるとノイズの影響を受けやすくなり，検出精度が悪化してしまう．
これに対し，サンプリング周波数の異なる疑似dq変換回路を3並列化することでノイズの影響を低減している．
本研究では，ノイズ成分が相互に低減されるサンプリング周波数の組み合わせを探索し，
PLL制御の制御特性の向上を目的とする．\\
　疑似dq変換は系統電圧に対しての近似微分であるため，系統に低次高調波が含まれた場合に大きく動脈
してしまう．この課題に対し，複素係数フィルタを用いて系統電圧の瞬時複素ベクトルを容易
に算出し，なおかつ，優れた高調波外乱に対する耐性を有したPLL制御が提案されている\cite{ref03}．
複素係数フィルタを用いたPLL制御の初期検討として，複素係数フィルタの制御特性について検証する．

